% !TeX root = ../../Book3.tex
% 纲要标题：### 17.4 应用与反例
% 纲要位置：计划纲要.md，第 2042 行。

\section{应用与反例}
\label{b3:sec:1704}

% 纲要内容：
%
% * 商环、理想及合冲模的计算
% * 不能对无限投射维数直接套公式
% * 为正则局部环的同调刻画准备工具
%
% ---
%

\subsection{商环、理想与合冲模的计算}
\label{b3:subsec:1704:01}

若 \(x_1,\ldots,x_c\in\mathfrak m\) 是交换 Noether 局部环 \(R\) 上的正则序列，则\cref{b3:thm:koszul-acyclic}给出商模的有限自由消解。因为各 \(x_i\in\mathfrak m\)，它极小，末项为 \(R\)，所以
\[
\operatorname{pd}_R R/(x_1,\ldots,x_c)=c,\qquad
\beta_i^R\bigl(R/(x_1,\ldots,x_c)\bigr)=\binom ci.
\]
Auslander–Buchsbaum 公式给出商模深度下降 \(c\)，与逐次除去非零因子的深度计算一致。

\ProseParagraphBreak
取 \(R=k[x,y]_{(x,y)}\)、\(I=(x,y)\)。上述 Koszul 消解去掉最后的 \(R\to k\)，得到
\[
0\longrightarrow R\xrightarrow{\binom{-y}{x}}R^2\longrightarrow I\longrightarrow0.
\]
因此 \(\operatorname{pd}_RI=1\)、\(\operatorname{depth}_RI=1\)。理想 \(I\) 是整环中的无挠模，却不是自由模；无挠只能排除单个非零元素的零化，不能保证深度与环一样大。

\begin{corollary}\label{b3:cor:syzygy-depth-finite-pd}
设 \((R,\mathfrak m)\) 为交换 Noether 局部环，\(M\ne0\) 为有限模且 \(r=\operatorname{pd}_RM<\infty\)。对 \(0\le j\le r\)，极小消解中的第 \(j\) 个合冲模非零，并满足
\[
\operatorname{pd}_R\Omega^jM=r-j,\qquad
\operatorname{depth}_R\Omega^jM=\operatorname{depth}_RM+j.
\]
\end{corollary}
\begin{proof}
截去极小消解的前 \(j\) 项得到 \(\Omega^jM\) 的极小消解，长度为 \(r-j\)，且末项仍非零。分别对 \(M\) 和 \(\Omega^jM\) 应用\cref{b3:thm:auslander-buchsbaum}，相减得到深度等式。
\end{proof}

\subsection{有限投射维数假设的必要性}
\label{b3:subsec:1704:02}

在 \(R=k[\epsilon]/(\epsilon^2)\) 上，\(\operatorname{depth}R=0\)，剩余域的深度也为零。然而\secref{b3:sec:1701}已算出 \(k\) 的每个 Betti 数都为一，故 \(\operatorname{pd}_Rk=\infty\)。不能把两个深度之差为零解释为 \(k\) 投射。

\ProseParagraphBreak
更有几何意味的例子是 \(R=k[x,y]_{(x,y)}/(xy)\) 与 \(M=R/(x)\)。\(R\) 中有
\(\operatorname{Ann}(x)=(y)\)、\(\operatorname{Ann}(y)=(x)\)：在多项式整环中约去 \(x\) 或 \(y\) 即得，局部化保持这些核。因此
\[
\cdots\xrightarrow{x}R\xrightarrow{y}R\xrightarrow{x}R
\longrightarrow M\longrightarrow0
\]
为极小自由消解，\(M\) 的投射维数无限。另一方面 \(M\cong k[y]_{(y)}\)，故 \(\operatorname{depth}_RM=1\)。\(R\) 的维数为一，而 \(x+y\) 不属于两个关联素理想 \((x),(y)\)，是非零因子，故 \(\operatorname{depth}R=1\)。深度相等仍不表示自由。

\subsection{正则局部环的同调刻画初步}
\label{b3:subsec:1704:03}

\begin{corollary}\label{b3:cor:finite-pd-maximal-depth-free}
设 \((R,\mathfrak m)\) 为交换 Noether 局部环，\(M\ne0\) 为有限 \(R\) 模。若 \(M\) 投射维数有限且 \(\operatorname{depth}_RM=\operatorname{depth}R\)，则 \(M\) 自由。特别地，深度为零的这种局部环上，每个具有有限投射维数的有限模都自由。
\end{corollary}
\begin{proof}
\cref{b3:thm:auslander-buchsbaum}使投射维数为零，因而模投射；局部有限投射模由\cref{b3:thm:finite-presentation-flat-projective}和\cref{b3:lem:finite-flat-local-free}自由。深度为零时，公式右边为零而左边两个整数均非负，故仍然如此。
\end{proof}

剩余域尤其敏感。若 \(\operatorname{pd}_Rk\) 有限，本章已经说明所有有限模的投射维数都有统一上界，且 \(\operatorname{pd}_Rk=\operatorname{depth}R\)。下一章要证明，这个纯同调条件还迫使 \(\dim R=\operatorname{edim}R\)。后一个等式不能仅从 Auslander–Buchsbaum 公式读出，需要另作关于极大理想的论证。
